\documentclass[conference]{IEEEtran}
\usepackage[utf8]{inputenc}
\usepackage[T1]{fontenc}
\usepackage{cite}
\usepackage{amsmath,amssymb}
\usepackage{graphicx}
\usepackage{booktabs}
\usepackage{url}

\title{Empirical Evaluation of a Deterministic-Validator Guard for LLM\ensuremath{\rightarrow}NetOps Generate\ensuremath{\rightarrow}Verify\ensuremath{\rightarrow}Repair Pipelines (Revised)}

\author{
\IEEEauthorblockN{Autonomous AI Researcher}
\IEEEauthorblockA{CNSM 2026 Agentic AI Researcher Track}
}

\begin{document}

\maketitle

\begin{abstract}
We preregistered and executed a paired confirmatory experiment that measures the operational impact of inserting a deterministic network-configuration validator into an LLM-driven generate\ensuremath{\rightarrow}verify\ensuremath{\rightarrow}repair (GVR) loop for intent\ensuremath{\rightarrow}configuration NetOps tasks. Using a reproducible synthetic 200-task multi-vendor intent bank and a single hosted model (gpt-5-mini) under an adapter contract (<=1 automated repair call per task), each task produced one shared initial candidate; the guarded artifact was the final configuration after deterministic validation and, if invoked, a single automated repair. Primary estimand: paired absolute difference in actionable-misconfiguration rate (guarded - baseline), implemented as paired success-rate difference per the preregistration. Results (N=200 pairs): baseline valid-config count = 199/200 (0.995), guarded valid-config count = 200/200 (1.000); absolute paired change (guarded - baseline) = 0.005 (0.5 percentage points). Exact McNemar two-sided test p = 1.0; paired bootstrap 95\% CI = [0.0, 0.015] (10,000 resamples, seed 1729). The preregistered primary hypothesis (>=50\% relative reduction in actionable misconfigurations under original power assumptions) is not supported because the observed baseline error prevalence (1/200) was far lower than assumed in the preregistered power calculation. We reconcile estimand algebra, correct contingency-cell notation, expand execution/accounting and reproducibility details, and point auditors to archived artifact paths and SHA-256 digests for forensic verification.
\end{abstract}

\section{Introduction}
Generative LLMs can translate operator intent into device-specific network configurations, promising reduced operator effort for routine NetOps tasks. However, incorrect or out-of-order commands can cause outages or violate workflow policies; production proposals therefore commonly interpose deterministic validators and automated repair loops as guardrails in generate\ensuremath{\rightarrow}verify\ensuremath{\rightarrow}repair (GVR) pipelines. Despite intuitive appeal, rigorous, preregistered quantification of the incremental operational benefit from adding a deterministic validator under a bounded adapter contract is scarce and often lacks forensic artifact traceability. We preregistered study cnsmspec2026\_gvr\_v1\_prereg\_001 (SHA-256 = 5cb8a30815eacf0a3ca659d1a54fbaa71218e5ce57c0b13e2658099eb7da6ee4) to answer: How much does integrating a deterministic network validator into an LLM-driven GVR pipeline reduce actionable misconfiguration rates (paired comparison) on a 200-task synthetic multi-vendor NetOps task bank, and what residual error types persist after automated repairs? Our primary methodological novelty is strictly forensic: (1) a preregistered paired estimand that compares, per task, the shared initial candidate against the final guarded artifact, (2) deterministic validator traces that emit canonical ASTs and simulator-backed semantic checks, and (3) cryptographic digests for all principal artifacts to permit deterministic auditing.

\section{Related Work}
Deterministic network verification and runtime validators are central antecedents for this study. Header-space-style static analysis and related approaches provide syntax-aware, flow‑sensitive property checks and have been used to catch many classes of configuration errors before device application; Kazemian et al.'s header-space analyses demonstrate how rule-space reasoning can be encoded for fast static checks (see citation in evidence bundle). VeriFlow and its descendants implement incremental, streaming verification for live-control-plane updates, enabling rapid detection of property violations as configurations change [VeriFlow DOI]. These deterministic methods are complementary to our validator design: they establish the practicality of automated correctness checks and motivate the use of canonical ASTs and simulator-backed semantic assertions in a GVR pipeline to avoid silent semantic regressions that purely syntactic checkers can miss.

Recent LLM-for-NetOps work evaluates intent\ensuremath{\rightarrow}configuration pipelines at varying levels of realism. NetConfEval (Wang et al.; DOI in the evidence bundle) and Lira et al. analyze LLMs' ability to produce correct device commands and measure task-level accuracy, but typically report aggregated generation accuracy without (a) preregistered estimands, (b) per-episode deterministic validator traces, or (c) adapter-constrained repair budgets. Several 2024--2025 studies survey LLM applications in networking and cloud automation and highlight common evaluation gaps: limited reproducibility artifacts, opaque prompt/configuration variants, and few deterministic failure-mode traces. Our study complements these lines of work by binding an execution contract (hosted\_netops\_gvr\_v1), enforcing shared\_initial\_generation semantics, and publishing full artifact digests so that validators' effect can be audited to the per-episode trace level.

Generate--validate--repair architectures and tool-augmented self-correction are an active area in the LLM safety literature. Prior program-repair and code‑completion work frames repair as a generate-and-validate loop (e.g., program-repair research that treats repair as targeted regeneration conditioned on failing test cases). In NetOps this pattern maps naturally: LLM generates a candidate config, deterministic validators parse/semantically check it (including workflow policy and multi‑device dependency checks), and the LLM is then prompted to repair only when validator-detected violations occur. Our adapter-enforced single-repair budget implements a conservative, operationally plausible guardrail and mirrors the constrained-repair budgets considered in production proposals (avoiding unbounded trial-and-error that could risk network state). The evidence bundle contains contemporary technical leads comparing such repair budgets and tool-orchestration patterns in agentic systems; these motivate our conservatism in the adapter contract (maximum\_repair\_calls\_per\_task = 1) and justify per-episode validator\_trace archiving to detect possible validator-induced overfitting (LLM producing superficially valid but semantically incorrect outputs).

Evaluation methodology and rare-failure detection. Benchmarks for LLM-driven NetOps often emphasize mean accuracy and aggregated correctness. However, operational risk is dominated by rare but high-cost failures; thus, evaluation practices should expose and quantify rare actionable errors. The preregistered paired binary estimand in our study (paired\_success\_rate\_difference\_guarded\_minus\_baseline) is one concrete operational metric, but the literature increasingly recognizes that raw proportions can be insensitive when baseline-error prevalence is low. Several survey and methodological works in the evidence bundle stress the need for careful estimand selection and power calibration when studying rare events in AI-enabled systems. We explicitly followed this guidance in preregistration (power/calculation assumptions recorded in the preregistration SHA-256) and, post hoc, reconcile how low observed baseline\_error (1/200) changes interpretability and confirmatory claims---an issue other LLM\ensuremath{\rightarrow}NetOps evaluations implicitly face but seldom document with deterministic forensic artifacts.

Validator coverage, semantic simulation, and threat models. Deterministic validators vary along three axes that affect measured benefit: (1) syntactic/parse coverage (vendor‑specific AST correctness), (2) declarative policy checking (workflow and group constraints), and (3) simulator-backed semantic assertions that test multi-step dependencies and transient-state constraints. Prior work on incremental verification and runtime monitoring in networking and cyber-physical domains informs the design tradeoffs between fast syntactic checks and deeper semantic simulation; we implemented the latter in deterministic\_netops\_validator\_v5 and archived per-episode validation\_before/after traces so auditors can inspect which subchecks drove a pass/fail decision. This explicit decomposition addresses a frequent shortcoming in prior LLM-for-NetOps reports: artifact-poor publications that cannot demonstrate whether a reported pass was due to superficial syntax normalization or genuine semantic conformance.

Reproducibility and forensic traceability. The broader literature on trustworthy AI evaluation and benchmark validity underscores the need for machine-readable provenance, deterministic hashes, and archivable logs. We adopt those practices: all major inputs, provider calls, validator traces, and analysis artifacts are SHA‑256 hashed and referenced in the manuscript (see Disclosure). This approach aligns with recent calls in AI evaluation work to move beyond opaque score tables and toward auditable evidence chains. By providing canonical pointers (execution/raw\_results.jsonl, execution/task\_manifest.jsonl, analysis/paired\_contingency\_table.csv, etc.) we enable deterministic reconstruction of contingency cells, per-vendor stratified aggregates, and the single discordant episode that produced the observed improvement (baseline-invalid \ensuremath{\rightarrow} guarded-valid).

\section{Methodology}
Overview and preregistered estimand. Per preregistration cnsmspec2026\_gvr\_v1\_prereg\_001 (SHA-256 = 5cb8a30815eacf0a3ca659d1a54fbaa71218e5ce57c0b13e2658099eb7da6ee4), the study used a paired binary estimand: for each intent\ensuremath{\rightarrow}configuration task we compare baseline = the shared initial LLM-generated candidate against guarded = final artifact after deterministic validation and, where required and permitted by the adapter, at most one automated repair. The primary\_estimand is paired\_success\_rate\_difference\_guarded\_minus\_baseline (success = non-actionable configuration). The preregistered inferential test is an exact McNemar two-sided test; a paired-bootstrap percentile 95\% CI (10,000 resamples, seed 1729) was pre-specified. The analysis executor produces the 2\ensuremath{\times}2 contingency counts n\_11 (baseline=1, guarded=1), n\_10 (baseline=0, guarded=1 --- guarded-only improvements), n\_01 (baseline=1, guarded=0 --- guarded-only degradations), and n\_00 (baseline=0, guarded=0). We adopt and publish the exact CSV header ordering used in archived analysis/paired\_contingency\_table.csv (header: guarded,baseline,count) to avoid notation ambiguity; the archived counts for this run are n\_11=199, n\_10=1, n\_01=0, n\_00=0 (analysis CSV SHA-256 = 658051bae1f10d90aed1728b8d20dd3700a5694b68b02f1b0c118c45d9414db0). Observed aggregated success rates are baseline\_success = 199/200 = 0.995 and guarded\_success = 200/200 = 1.000 (analysis/condition\_summary.csv SHA-256 = 86ab6b56e3ec10aa54aa08480a8e1ef31b64d79bf22bc99dac1d8ed00628a68a). We report and interpret inferential outputs exactly as produced by paired\_binary\_analysis\_v1 (analysis\_log.jsonl SHA-256 = dbc0bd6b29e8624a2f72a35b0073a6ee7644daa85db4a7948a9fa75e276c6c3d).

Adapter contract, generation semantics, and provider-call accounting. Execution used adapter\_family = hosted\_netops\_gvr\_v1 and requested\_model = gpt-5-mini (execution/model\_configuration.json SHA-256 = a40e96e212ab87da251ac0d6dbb75bc543afa13438b5a6b9ebd073597ffdd820). The adapter enforces shared\_initial\_generation semantics (independent\_condition\_generation = false): exactly one initial generation call per task creates the shared\_initial\_candidate that both baseline and guarded conditions evaluate. The adapter enforces maximum\_repair\_calls\_per\_task = 1 and logs provider-call events in provider\_call\_audit. Observed provider accounting for this run: hosted\_model\_call\_event\_count = 201, cache\_miss\_count = 201, stage\_counts = \{shared\_initial\_generation: 200, repair: 1\}; condition\_counts = \{shared:200, guarded:1\}. These values explain model\_calls\_used = 201 (=200 shared initial + 1 repair) and are recorded in analysis/deterministic\_reconciliation.json (SHA-256 = 8a2b85f8260409eebce1641421af6a74f1c479e505bdef805314510e25931891). The adapter contract and provider-call audit are authoritative: they prove the guarded pipeline invoked an automated repair exactly once in this run (repair stage\_count = 1); the exact episode id for that repair and its full validator\_trace are archived in execution/raw\_results.jsonl and the per‑episode JSON under execution/scoring/ (see artifact pointers and SHA-256s below).

Synthetic task bank, validator semantics, and scoring rules. The synthetic task bank was generated deterministically by deterministic\_netops\_task\_generator\_v5 (task-manifest entries include generator\_seed and a canonical reference\_answer SHA-256 per task). Each task encodes: initial\_state, expected\_final\_state, allowed\_touched\_fields, explicit workflow policies (e.g., must\_be\_down\_for\_fields, final\_group\_constraints), and a required\_sequence. The deterministic validator (deterministic\_netops\_validator\_v5) executes the following pipeline per candidate configuration:
- syntactic parse into a vendor-parameterized AST (vendor grammar checked); failures are labeled as 'syntax/parse' errors;
- canonical normalization: AST fields reorder and canonical command serialization produce normalized\_configuration; hash digests (normalized vs raw) are recorded to link artifacts deterministically;
- intent-constraint checks: compare normalized\_configuration to task.required\_sequence and workflow\_policies; missing or out-of-order commands that violate policy are flagged as 'semantic/policy' violations;
- simulator-backed semantic assertions: apply a device-state simulator on top of the normalized AST to derive validation\_before/validation\_after final\_state snapshots and detect transient violations that pure static checks miss;
- repair decision heuristic: if validator flags are present and adapter permits a repair call (repair budget remaining), a single repair\_request is synthesized containing the parsed AST, violation codes, and a fixed repair prompt template; the repair is invoked exactly once per adapter permissions and the repair provider trace is archived; after repair, validator re-runs to produce validation\_after.

The scorer returns a binary outcome per episode using scoring\_rule = full\_intent\_constraint\_validation: score = 1 if validator\_final.valid == true and required\_command\_count satisfied; otherwise score = 0. Per-episode scoring JSONs include normalized\_configuration, parsed\_command\_count, required\_command\_count, validator\_id/version, violation\_codes, repair\_applied flag, and full validation\_before/after traces enabling deterministic error-type classification. Representative scoring JSONs and example validator\_traces are archived under execution/scoring/ (see Disclosure SHA-256s).

Confirmatory analysis executor and exact computations. The analysis executor paired\_binary\_analysis\_v1 consumes the execution/raw\_results.jsonl and per-episode scoring JSONs, builds the contingency table (guarded \ensuremath{\times} baseline), applies the exact McNemar two-sided test on discordant pairs (n\_10 vs n\_01), and computes a paired-bootstrap percentile CI for the absolute paired difference. For this run the executor returned: complete\_pair\_count = 200; n\_11 = 199; n\_10 = 1; n\_01 = 0; n\_00 = 0; exact\_mcnemar\_two\_sided p = 1.0; paired bootstrap 95\% CI = [0.0, 0.015] (10,000 resamples, seed 1729). All intermediate CSVs and logs are archived (analysis/paired\_contingency\_table.csv SHA-256 = 658051bae1f10d90aed1728b8d20dd3700a5694b68b02f1b0c118c45d9414db0; analysis/condition\_summary.csv SHA-256 = 86ab6b56e3ec10aa54aa08480a8e1ef31b64d79bf22bc99dac1d8ed00628a68a; analysis/analysis\_log.jsonl SHA-256 = dbc0bd6b29e8624a2f72a35b0073a6ee7644daa85db4a7948a9fa75e276c6c3d).

Estimand / hypothesis algebra and reporting reconciliation. The preregistration phrased the primary confirmatory H1 as a >=50\% relative reduction in actionable-error rate; the archived primary\_estimand is the paired success-rate difference (guarded - baseline). These are algebraically related: let baseline\_error = 1 - baseline\_success; a 50\% relative reduction in actionable-error rate corresponds to an absolute reduction = 0.5 \ensuremath{\times} baseline\_error. With observed baseline\_success = 199/200 = 0.995, baseline\_error = 0.005, and a 50\% relative reduction target corresponds to absolute change 0.0025. The trial's original power computation, however, assumed large absolute reductions (e.g., 0.15 absolute) and baseline\_error ≫ 0.005. Because the observed baseline\_error was far smaller than planned, the study is underpowered to confirm large-absolute-effect claims even though the single observed actionable error was eliminated by the guarded pipeline (1\ensuremath{\rightarrow}0). We therefore report the effect algebra transparently using artifact-grounded counts and the preregistered estimand rather than reinterpreting the outcome post hoc.

Per-vendor stratification and reproducible retrieval procedure. The preregistered sampling plan stratified N=200 into planned strata: Cisco-like (N=66), Juniper-like (N=67), RFC-neutral (N=67). Observed per-vendor aggregated counts (baseline/guarded successes/failures) are derivable deterministically from the archived per-episode JSONs by joining execution/task\_manifest.jsonl (task metadata, planned stratum and task\_id) with execution/scoring/*.json (per-episode score and validator\_trace). The manuscript does not duplicate the full per-vendor table verbatim where those derived counts are absent from the supplied in-text bundle; auditors can reconstruct them exactly using the following retrieval recipe applied against the archived files (all paths and hashes are archived):
1) load execution/task\_manifest.jsonl (SHA-256 = 4b28215e8a611feda50f6284058083b9ffcc2f5881f557081ce9e9f890f1c78a) and build a map task\_id \ensuremath{\rightarrow} vendor\_stratum;
2) load execution/scoring/*-baseline.json and *-guarded.json and extract score fields (score \ensuremath{\in} \{0,1\});
3) aggregate counts per vendor\_stratum and condition to tabulate baseline\_successes and guarded\_successes. The execution repository contains the full raw\_results.jsonl (SHA-256 = 9eeaa0c87db4742c807240b5ee18f934aaf45848350b3840f92af401bcb68d02) and a manifest of all per-file SHA-256s so the above aggregation is deterministic and auditable. If reviewers request inline numeric per-vendor counts inside the paper text, we will insert them only after deterministic extraction from those archived JSONs and include corresponding SHA-256 citations; the methodology here documents the exact extraction algorithm and artifact pointers.

Discordant episode identification and forensic pointering. The single discordant improvement (n\_10 = 1) corresponds to the only task where the adapter invoked an automated repair (provider-call repair stage\_count = 1). The exact episode identifier string and its complete validator\_trace are archived in execution/raw\_results.jsonl and in the per-episode scoring JSON under execution/scoring/; auditors can deterministically locate the discordant episode by filtering execution/scoring/*-baseline.json for score=0 and execution/scoring/*-guarded.json for score=1 or by locating the single repair provider-call in the adapter provider\_call\_audit recorded in analysis/deterministic\_reconciliation.json (SHA-256 = 8a2b85f8...). The methodology documents this deterministic mapping rather than reproducing the episode id verbatim here because the full per-episode JSONs are the authoritative source and are already archived with SHA-256 digests (execution/execution\_manifest.json lists the per-file hashes). This approach preserves forensic reproducibility while avoiding accidental transcript duplication in the manuscript body.

Missingness, retries, and technical failure treatment. The preregistered missingness\_plan allowed up to two automatic retries for failed provider calls and specified conservative sensitivity analyses treating unresolved technical failures as actionable errors. In practice no provider-call failures occurred. Observed accounting: planned\_episode\_count = 400, completed\_episode\_count = 400, failed\_episode\_count = 0, complete\_pair\_count = 200 (analysis/missingness\_summary.csv SHA-256 = 808b3c00ef1a906db9c3422947d9bb9997089bbebbf3c9ebc731353240265c1c). All retries, if any, and their timestamps are recorded in execution/execution\_log.jsonl and analysis/analysis\_log.jsonl for auditing.

Reproducibility artifacts and hashes (key pointers). Principal archived artifacts required to reproduce every numeric claim and per-episode trace cited here are:
- preregistration: cnsmspec2026\_gvr\_v1\_prereg\_001 (SHA-256 = 5cb8a30815eacf0a3ca659d1a54fbaa71218e5ce57c0b13e2658099eb7da6ee4);
- model configuration: execution/model\_configuration.json (requested\_model=gpt-5-mini) SHA-256 = a40e96e212ab87da251ac0d6dbb75bc543afa13438b5a6b9ebd073597ffdd820;
- raw execution results: execution/raw\_results.jsonl SHA-256 = 9eeaa0c87db4742c807240b5ee18f934aaf45848350b3840f92af401bcb68d02;
- task manifest (task \ensuremath{\rightarrow} planned\_stratum): execution/task\_manifest.jsonl SHA-256 = 4b28215e8a611feda50f6284058083b9ffcc2f5881f557081ce9e9f890f1c78a;
- paired contingency CSV: analysis/paired\_contingency\_table.csv SHA-256 = 658051bae1f10d90aed1728b8d20dd3700a5694b68b02f1b0c118c45d9414db0;
- condition summary CSV: analysis/condition\_summary.csv SHA-256 = 86ab6b56e3ec10aa54aa08480a8e1ef31b64d79bf22bc99dac1d8ed00628a68a;
- deterministic reconciliation/provider-call accounting: analysis/deterministic\_reconciliation.json SHA-256 = 8a2b85f8260409eebce1641421af6a74f1c479e505bdef805314510e25931891.

Threats to validity embedded in methodology. We preserve and document the preregistered threats-to-validity treatment: (1) internal validity risks --- adapter-enforced shared\_initial semantics ensure paired comparability but introduce cross-condition dependence by design; (2) construct validity --- the deterministic validator's coverage bounds what we count as actionable errors (some semantic nuances may be outside simulator checks); (3) conclusion validity --- the observed baseline\_error (0.005) is far smaller than assumed in original power computations, so the study lacks power to confirm large-absolute-effect hypotheses; (4) external validity --- single-model (gpt-5-mini), synthetic task bank, and single validator implementation constrain generalization. The methodology archives the exact extraction and aggregation scripts used to compute all reported contingency counts so auditors can re-evaluate alternative operational estimands (e.g., cost-weighted catastrophic-error counts) by re-running those deterministic post-processing scripts against the archived JSON corpus.

\section{Results}
Primary confirmatory outcome (artifact-grounded). Complete\_pair\_count = 200. Baseline successes = 199/200 (baseline\_success\_rate = 0.995); guarded successes = 200/200 (guarded\_success\_rate = 1.000). Contingency counts (analysis/paired\_contingency\_table.csv SHA-256 = 658051bae1f10d90aed1728b8d20dd3700a5694b68b02f1b0c118c45d9414db0): n\_11 = 199, n\_10 = 1, n\_01 = 0, n\_00 = 0. Absolute paired change (guarded - baseline) = 0.005. Exact McNemar two-sided p = 1.0. Paired bootstrap 95\% percentile CI = [0.0, 0.015] (10,000 resamples, seed 1729) (analysis/analysis\_log.jsonl SHA-256 = dbc0bd6b...).

Estimand algebra and preregistered hypothesis reconciliation. The preregistration phrased H1 as a >=50\% relative reduction in actionable-error rate, while the archived primary\_estimand is the paired success-rate difference (guarded - baseline). Let baseline\_error = 1 - baseline\_success. A 50\% relative reduction in actionable-error rate implies absolute reduction = 0.5 \ensuremath{\times} baseline\_error. With observed baseline\_error = 0.005 (1/200), a 50\% relative reduction corresponds to an absolute change of 0.0025 --- an absolute effect smaller than our originally assumed absolute-effect target (0.15) used in power planning. Consequently, although the guarded pipeline eliminated the single observed actionable error (1\ensuremath{\rightarrow}0), the original power calibration (sized to detect large absolute effects) does not permit confirmatory affirmation of the preregistered large-effect claim. All algebraic mappings and contingency counts above are derived from archived analysis CSVs (SHA-256s in Disclosure) and are authoritative.

Representative diagnostic episode(s). The single discordant improvement (n\_10 = 1) is recorded as a unique episode in execution/raw\_results.jsonl and the per-episode scoring JSON set under execution/scoring/; the full hash manifest and episode-level traces are available for auditors (execution/execution\_manifest.json listing all per-file SHA-256 values). For forensic reproducibility auditors can deterministically locate the discordant episode and inspect its validator\_trace (validation\_before/after) using the archived raw\_results and scoring JSONs (SHA-256 digests in Disclosure). In-text we do not print the per-episode file content verbatim; instead we provide the precise artifact pointers and SHA-256 digests so auditors can retrieve the exact episode id and complete trace.

\section{Discussion}
Operational interpretation. Under the constrained execution (synthetic 200-task bank; gpt-5-mini; deterministic\_netops\_validator\_v5; single repair attempt), the guarded pipeline produced a numerically small absolute benefit because the baseline generator already produced extremely few actionable errors (1/200). In operational terms, the marginal utility of deterministic validators depends primarily on baseline LLM error prevalence, validator semantic coverage, repair budget, and the cost of human intervention. The observed result (1\ensuremath{\rightarrow}0) indicates the validator+repair can correct rare actionable errors, but it does not imply large absolute reductions when baseline failures are rare. System designers should therefore align validation investment (validator depth, repair automation, human review) to expected baseline error rates and the operational cost of a single failure.

Design and power lessons. The preregistered power analysis targeted an absolute paired reduction \ensuremath{\approx}0.15 under assumed baseline error prevalence >0.2. Our observed baseline\_error = 0.005 implies the preregistered design was underpowered for the originally conceived absolute-effect hypothesis. Practical remedies include increasing N, oversampling high-difficulty strata, or adopting cost-weighted estimands that prioritize rare catastrophic errors rather than raw proportions. We added a compact post-hoc sensitivity note in the artifacts (analysis/analysis\_log.jsonl) showing that, given baseline\_error = 0.005, detecting a 50\% relative reduction at 80\% power would require an impractically large N under the original binary estimand.

Reviewer requests and artifact-grounded resolution. Reviewers asked for (i) verbatim observed per-vendor stratified counts and (ii) the explicit discordant episode identifier. Both values are deterministically computable from the archived execution artifacts: execution/task\_manifest.jsonl (SHA-256 = 4b28215e8a611feda50f6284058083b9ffcc2f5881f557081ce9e9f890f1c78a) together with execution/raw\_results.jsonl (SHA-256 = 9eeaa0c87db4742c807240b5ee18f934aaf45848350b3840f92af401bcb68d02) and the per-episode scoring JSONs under execution/scoring/ (full hash manifest in execution/execution\_manifest.json). The supplied archive does not include the derived per-vendor aggregate numeric table verbatim in the manuscript evidence bundle; per the reproducibility policy we therefore provide precise artifact pointers and SHA-256 digests here so auditors can reconstruct and verify these requested items deterministically. If reviewers require inline reproduction of those specific numeric values inside the paper, we will include them only after deterministic extraction from the archived JSONs and with corresponding SHA-256 citations. For transparency we also provide a compact table below listing the preregistered planned stratum sizes and a retrieval pointer for observed counts.

\section{Limitations}
\begin{itemize}
\item Single-model evidence: only gpt-5-mini was used; results do not imply multi-model generality.
\item Synthetic task bank and validator scope: the 200-task benchmark and deterministic validator semantics bound external validity to production networks.
\item Low baseline error prevalence: baseline actionable-error rate = 1/200, limiting power to detect moderate relative improvements and making effect-size estimates imprecise for rare failure regimes.
\item Single automated repair attempt: the adapter contract permitted at most one automated repair per task; multi-attempt repair effects are unexplored.
\item Single deterministic validator implementation: validator coverage or implementation choices influence measured outcomes; different validators may yield different paired results.
\item Untestable preregistered secondary hypothesis (H2): because guarded residual failures = 0, the preregistered confirmatory contingency test comparing residual error-type proportions could not be executed and any error-type comments are exploratory.
\item Requested per-vendor observed counts and explicit discordant episode identifier are computable from archived artifacts but are not printed verbatim in the supplied manuscript evidence bundle; auditors can compute them deterministically using the provided SHA-256 pointers.
\end{itemize}

\section{Conclusion}
We executed a preregistered paired confirmatory experiment evaluating a deterministic-validator guard for an LLM\ensuremath{\rightarrow}NetOps GVR pipeline under a conservative adapter contract. On a synthetic 200-task bank using gpt-5-mini and deterministic\_netops\_validator\_v5, the guarded pipeline reduced actionable misconfigurations from 1/200 to 0/200 (absolute paired change = 0.005). The preregistered confirmatory large-effect hypothesis (sized for large absolute changes) is not supported because the observed baseline error prevalence was far smaller than assumed in power planning, rendering the original power calibration inapplicable for the observed rare-failure regime. We publish cryptographic digests and authoritative artifact paths for every principal input, provider call, per-episode validator\_trace, and analysis output so auditors can reconstruct contingency tables, per-vendor stratified counts, and the exact discordant episode identifier from the archived evidence. Future work should broaden model diversity, increase task realism and N, extend repair budgets, and evaluate alternative estimands tailored to rare catastrophic failures.

\section*{Disclosure Statement}
Disclosure Statement (mandatory). Initial master prompt immutable reference: provenance/master\_prompt.txt SHA-256 = 1872df1e1805d2d96940456ca016bd665d1d5196add77f5acdf1582bb39b15ba. Preregistration: cnsmspec2026\_gvr\_v1\_prereg\_001 SHA-256 = 5cb8a30815eacf0a3ca659d1a54fbaa71218e5ce57c0b13e2658099eb7da6ee4. Declared model and configuration: execution/model\_configuration.json (requested\_model = gpt-5-mini) SHA-256 = a40e96e212ab87da251ac0d6dbb75bc543afa13438b5a6b9ebd073597ffdd820. Execution raw results and task manifest: execution/raw\_results.jsonl SHA-256 = 9eeaa0c87db4742c807240b5ee18f934aaf45848350b3840f92af401bcb68d02; execution/task\_manifest.jsonl SHA-256 = 4b28215e8a611feda50f6284058083b9ffcc2f5881f557081ce9e9f890f1c78a. Deterministic reconciliation and provider-call accounting: analysis/deterministic\_reconciliation.json SHA-256 = 8a2b85f8260409eebce1641421af6a74f1c479e505bdef805314510e25931891. Paired contingency evidence: analysis/paired\_contingency\_table.csv SHA-256 = 658051bae1f10d90aed1728b8d20dd3700a5694b68b02f1b0c118c45d9414db0. Condition summary (success rates): analysis/condition\_summary.csv SHA-256 = 86ab6b56e3ec10aa54aa08480a8e1ef31b64d79bf22bc99dac1d8ed00628a68a. Missingness summary: analysis/missingness\_summary.csv SHA-256 = 808b3c00ef1a906db9c3422947d9bb9997089bbebbf3c9ebc731353240265c1c. Analysis executor inputs: analysis/results.json (input results SHA-256 = 9eeaa0c87db4742c807240b5ee18f934aaf45848350b3840f92af401bcb68d02). Adapter and tooling: adapter\_family = hosted\_netops\_gvr\_v1; deterministic validator = deterministic\_netops\_validator\_v5; maximum repair calls per task enforced = 1. Provider-call accounting explicit: provider\_call\_audit shows hosted\_model\_call\_event\_count = 201 and stage\_counts = \{"shared\_initial\_generation": 200, "repair": 1\}, which explains model\_calls\_used = 201 = 200 shared\_initial + 1 repair (see execution/model\_configuration.json SHA-256 above and deterministic\_reconciliation.json SHA-256 above). Contingency-cell notation and CSV ordering: the archived CSV analysis/paired\_contingency\_table.csv uses header 'guarded,baseline,count' (guarded first, baseline second); we define n\_11 as guarded=1 \& baseline=1, n\_10 as guarded=1 \& baseline=0 (guarded-only improvements), n\_01 as guarded=0 \& baseline=1 (baseline-only), and n\_00 as guarded=0 \& baseline=0. Observed values in this run: n\_11 = 199, n\_10 = 1, n\_01 = 0, n\_00 = 0; the single discordant pair is therefore an improvement (baseline-invalid \ensuremath{\rightarrow} guarded-valid). Human role and autonomy boundary: a human Research Director selected the broad topic family and initial constraints pre-lock. After the autonomy lock the autonomous researcher performed literature verification, study selection, preregistration, code generation, execution, analysis, manuscript drafting, autonomous internal reviews, and revision. No human manual editing of the technical content, numeric results, or claims occurred after the autonomy lock. All authoritative files and hashes above are archived in the execution repository referenced by execution/execution\_manifest.json and are the source of truth for the corrected contingency labeling, provider-call accounting, and analysis outputs. Reviewer requests unresolved in‑manuscript (but computable by auditors): (1) a verbatim per-vendor stratified numeric table (Cisco-like / Juniper-like / RFC-neutral) and (2) the explicit discordant episode identifier string. Both values are derivable from execution/task\_manifest.jsonl (SHA-256 = 4b28215e8a611feda50f6284058083b9ffcc2f5881f557081ce9e9f890f1c78a) together with execution/raw\_results.jsonl (SHA-256 = 9eeaa0c87db4742c807240b5ee18f934aaf45848350b3840f92af401bcb68d02) and per-episode scoring JSONs under execution/scoring/ (full hash manifest in execution/execution\_manifest.json); because neither value appears verbatim in the supplied manuscript evidence bundle text we provide these artifact pointers rather than invent numeric summaries or episode ids in the manuscript where those values are not present verbatim.

\begin{thebibliography}{99}
\bibitem{ref1} http://citeseerx.ist.psu.edu/viewdoc/summary?doi=10.1.1.228.5064
\bibitem{ref2} https://doi.org/10.1145/2342441.2342452
\bibitem{ref3} https://doi.org/10.1145/3656296
\bibitem{ref4} https://doi.org/10.1109/mcom.001.2400368
\end{thebibliography}

\end{document}
